% !TEX root = paper.tex


\section{\verus by example} \label{sec:design}

This section introduces the basic features of \verus
by walking through a simple example that computes Fibonacci numbers,
shown in \autoref{listing:fibonacci}.
The example consists of a set of functions written in Rust.
Each function is annotated with an attribute,
using Rust's \verb`#[...]` attribute syntax,
to indicate whether the function is executable code (\verb`#[exec]`),
proof code (\verb`#[proof]`),
or specification code (\verb`#[spec]`).
We refer to \verb`exec`, \verb`proof`, and \verb`spec` as modes;
\autoref{fig:modes} summarizes the properties of these three modes.

\begin{figure}

\noindent\begin{minipage}{.5\textwidth}

\begin{lstlisting}[language=Verus,style=VerusLineNos]
#[spec] fn fibo(n: nat) -> nat {
    decreases(n); %*\clnum\label{clnum:fibo:fibo-decreases}*
    if n == 0 { 0 }
    else if n == 1 { 1 }
    else { fibo(n - 2) + fibo(n - 1) }
}

#[proof] fn lemma_fibo_is_monotonic(i:nat, j:nat) { 
    requires(i <= j); %*\clnum\label{clnum:fibo:requires-clause}*       %* \codemarker{spec}*
    ensures(fibo(i) <= fibo(j)); %*\clnum\label{clnum:fibo:ensures-clause}*       %* \codemarker{spec}*
    decreases(j - i); %*\clnum\label{clnum:fibo:decreases-clause}*       %* \codemarker{spec}*

    if (i < 2 && j < 2) || i == j {
    } else if i == j - 1 {
        reveal_with_fuel(fibo, 2); %*\clnum\label{clnum:fibo:reveal}*
        lemma_fibo_is_monotonic(i, j - 1); %*\clnum\label{clnum:fibo:lemma-monotonic-induct}*
    } else {
        lemma_fibo_is_monotonic(i, j - 1);
        lemma_fibo_is_monotonic(i, j - 2);
    }
}

#[spec] fn fibo_fits_u64(n: nat) -> bool {
    fibo(n) <= u64::MAX
} 
\end{lstlisting}
\end{minipage}\hfill
\noindent\begin{minipage}{.5\textwidth}

\begin{lstlisting}[language=Verus,style=VerusLineNos,firstnumber=26]
#[exec] fn fibo_impl(n: u64) -> u64 {
    requires(fibo_fits_u64(n)); %*\clnum\label{clnum:fibo:requires-clause-impl}*       %* \codemarker{spec}* 
    ensures(|result: u64| result == fibo(n));%*\clnum\label{clnum:fibo:ensures-clause-impl}*       %* \codemarker{spec}* 

    if n == 0 { return 0; } 
    let mut prev: u64 = 0;
    let mut cur: u64 = 1;
    let mut i: u64 = 1;
    while i < n {
        invariant([ %*\clnum\label{clnum:fibo:invariant-clause}*       %* \codemarker{spec}* 
            0 < i && i <= n,
            fibo_fits_u64(n as nat),
            fibo_fits_u64(i as nat),
            cur == fibo(i),
            prev == fibo(i as nat - 1),
        ]);
        let new_cur = cur + prev;
        prev = cur;
        cur = new_cur;
        assert(prev == fibo(i as nat)); %*\clnum\label{clnum:fibo:fibo-impl-assert}*       %* \codemarker{proof}*
        i = i + 1;
        lemma_fibo_is_monotonic(i, n); %*\clnum\label{clnum:fibo:fibo-impl-lemma}*       %* \codemarker{proof}*
    }
    cur
}
\end{lstlisting}
\end{minipage}
\caption{A proof of correctness of a function computing the n-th
Fibonacci number. We use circled letters, similar to \protect\circled{Z},
to mark points of interest in the code. The markers \protect\inlinecodemarker{spec} and
\protect\inlinecodemarker{proof} indicate specification and proof mode code respectively
when it differs from the mode of the function.}
\label{listing:fibonacci}
\end{figure}

In \verus, specifications and proofs are simply Rust code,
parsed with Rust's parser and checked with Rust's type checker.
This avoids the need for systems programmers to learn a separate verification language,
making verification more accessible and convenient.
It also allows specifications and proofs to take advantage of Rust's features,
such as recursive functions, arithmetic, algebraic datatypes,
pattern matching, modules, closures, traits, etc.
For soundness's sake, \verus places some limits
on the features that specifications and proofs can use.
In particular, specifications must be deterministic,
and recursive \verb`spec` functions and recursive \verb`proof` functions must terminate.

The \verus tool, which extends the Rust compiler,
erases all ghost code (all specifications and proofs) before the code is compiled to machine code.
In the example, \verb`lemma_fibo_is_monotonic`, \verb`fibo`, and \verb`fibo_fits_u64`
are all erased before compilation.
Furthermore, the executable function \verb`fibo_impl` contains small bits of specification and proof
inside its body
(\clref{clnum:fibo:requires-clause-impl},
\clref{clnum:fibo:ensures-clause-impl},
\clref{clnum:fibo:invariant-clause},
\clref{clnum:fibo:fibo-impl-assert},
\clref{clnum:fibo:fibo-impl-lemma}),
and this ghost code is also erased.

\verus encodes preconditions and postconditions as calls to \verus library functions
named \verb`requires` and \verb`ensures`.
Postconditions may refer to the return value;
for this, the ensures function accepts a Rust closure that declares a name for the return value.
(In Rust, first-class functions are called closures and have the syntax ``\verb`|...parameters...| body`''.)
The example uses a postcondition to prove that the executable function \verb`fibo_impl`
computes the same result as the mathematical definition of the $n$\textsuperscript{th} Fibonacci number in the
\verb`fibo` function: the \verb`ensures` clause~\clref{clnum:fibo:ensures-clause-impl}
establishes this postcondition for \verb`fibo_impl`. Because the return value is a bounded 64-bit unsigned
integer, \verb`fibo_impl` can only accept a parameter \verb`n` such that the n-th Fibonacci number fits
in the type of the return value: this is established by the \verb`requires` clause~\clref{clnum:fibo:requires-clause-impl}.
Note that \verus extends Rust's type system with two new integer types,
\verb`int` (mathematical integers $\mathbb{Z}$), and \verb`nat` (natural numbers $\mathbb{N}$),
so that specifications and proofs can talk about arbitrary integers.
Executable code, however, is limited to Rust's finite-width integer types like \verb`u64`,
since \verb`int` and \verb`nat` aren't compilable to machine code.

To help prove the postcondition,
the \verb`fibo_impl` function uses a loop invariant \clref{clnum:fibo:invariant-clause}
containing a list of clauses that must be true before and after each loop iteration.
Given preconditions, postconditions, and loop invariants,
\verus uses standard weakest precondition reasoning~\cite{DBLP:journals/cacm/Dijkstra75}
to generate a verification condition for \verb`fibo_impl`.
It then sends this verification condition to the Z3 SMT solver~\cite{z3}.

In many cases, the SMT solver can prove the verification condition completely automatically.
In other cases, the proof may require reasoning beyond the SMT solver's abilities.
For example, to prove the absence of 64-bit integer overflow,
\verb`fibo_impl` relies on the Fibonacci sequence being monotonic,
which requires an inductive proof that the SMT solver cannot generate automatically.
Instead, the programmer supplies an inductive proof in the form of a recursive \verb`proof` function
(e.g., \verb`lemma_fibo_is_monotonic`).
The programmer can also add explicit assertions\clref{clnum:fibo:fibo-impl-assert} that
serve as hints to the SMT solver.
This style of SMT-based verification with programmer-supplied lemmas and hints
is similar to other verification systems like Boogie~\cite{boogie}, Dafny, and \fstar.

To improve verification performance,
\verus strives to keep the verification condition encoding lightweight,
so that the SMT encoding of specifications is not much larger than the original
specifications written in \verus code.
In particular, calls to \verb`spec` functions are translated directly into calls to SMT functions,
with no additional overhead.
For this reason, \verus \verb`spec` functions are total functions that do not have preconditions and postconditions.
This design is similar to Boogie,
though it differs from Dafny and \fstar.

This design choice has a downside, since
precondition failures can provide the developer with early feedback
to find errors in specification functions and in how they are used.
In order to restore that feedback, \verus introduces \verb`recommends` clauses: soft preconditions
for \verb`spec` functions, which Verus only considers when there is a verification error.
At that point, it performs a separate check for soft preconditions of
\verb`spec` functions that are mentioned in the context of the failure,
and reports failures as warnings for the developer.

Recursive \verb`spec` functions and recursive \verb`proof` functions are valid only if they terminate on all inputs
(otherwise, they could encode unsound circular reasoning).
\verus requires that all such functions contain a \verb`decreases` clause~\clref{clnum:fibo:fibo-decreases}~\clref{clnum:fibo:decreases-clause}
and each recursive call must decrease the expression in the clause. The recursive definition of
the $n$\textsuperscript{th} Fibonacci number~\clref{clnum:fibo:fibo-decreases} in \autoref{listing:fibonacci}
is legal because both recursive calls decrease the expression \verb`n`.
(\verus also imposes positivity restrictions on recursive type definitions to prevent
nontermination, as discussed in \autoref{sec:formalization}.)
The SMT solver may need to unfold definitions of a recursive \verb`spec` function.
As in Dafny and \fstar, \verus uses an integer ``fuel'' to control the number of unfoldings.
The \verb`reveal_with_fuel`~\clref{clnum:fibo:reveal} function controls the fuel level.

\begin{figure}
  \centering
  \small
  \begin{tabular}{r|l|l|l}
    \hline
    ~ & specification mode & proof mode & executable mode \\ \hline
    compiled or ghost & ghost & ghost & compiled \\
    code style & purely functional & mutation allowed & mutation allowed \\
    linearity \& borrowing checking & not checked & checked & checked \\
    can call specification functions & yes & yes & yes \\
    can call proof functions & no & yes & yes \\
    can call executable functions & no & no & yes \\
    % body visibility & may be visible & never visible & never visible \\
    % body & body optional & body mandatory & body mandatory \\
    determinism & deterministic & nondeterministic & nondeterministic \\
    termination & must terminate & must terminate & nontermination ok \\
    preconditions/postconditions & none & requires/ensures & requires/ensures \\
    \hline
  \end{tabular}
  \caption{Summary of \verus' modes and their properties.}
  \label{fig:modes}
\end{figure}

\subsection{Linearity, Borrowing, Spec Variables, and Proof Variables} \label{sec:design:linearity}

Rust types are linear by default:
unless a type implements the Rust \verb`Copy` trait,
values of the type can only be moved from one variable to another, not copied.
For example, the Rust \verb`Vec<T>` type for vectors is linear.
The following code is illegal in Rust because it attempts to duplicate a \verb`Vec<u64>` value,
returning both copies of the value in a pair:

\begin{lstlisting}[language=Verus]
#[exec] fn f(v: Vec<u64>) -> (Vec<u64>, Vec<u64>) {
    let v1 = v;
    let v2 = v; // illegal, tries to duplicate v
    (v1, v2)
}
\end{lstlisting}

On the other hand, Rust code can duplicate immutable references to values,
as long as the scope of the references is limited.
In Rust terminology, a reference of type \verb`&T` temporarily {\it borrows} from
an owned value of type \verb`T`.
During the borrowing, the original owned value of type \verb`T` is inaccessible.
When the references go out of scope, the original owned value becomes accessible again.
Rust code can also borrow a mutable reference of type \verb`&mut T`;
in contrast to immutable references, mutable references cannot be duplicated.
Rust enforces the property that a value cannot be borrowed both immutably and mutably simultaneously.

Rust contains a sophisticated ``borrow checker'' that checks linearity and borrowing.
Similar to Creusot~\cite{denis:hal-03737878}, \verus trusts the results of Rust's borrow checker,
and does not attempt to recheck these results in the SMT solver,
since this would just slow down the SMT solving.
Because of this, \verus can rely on the properties of borrowing in its SMT encoding.
For example, \verus encodes immutably borrowed references \verb`&T` and owned heap pointers
(\verb`Box<T>`, \verb`Rc<T>`, and \verb`Arc<T>`) simply as values of type \verb`T`,
not as pointers to locations.

\verus specifications, in contrast to ordinary Rust code, are not checked for linearity and borrowing;
specifications can freely copy any value of any type.
This allows specifications to freely talk about linear values,
potentially mentioning a single linear variable multiple times in a precondition or postcondition,
for example.
\verus code can also store nonlinear copies of linear variables inside {\it spec variables},
declared with the attribute \verb`#[spec]`:

\begin{lstlisting}[language=Verus]
#[exec] fn f(v: Vec<u64>) {
    #[spec] let v1 = v; // copies v into spec variable v1
    #[spec] let v2 = v; // copies v into spec variable v2
    assert(v1.len() == v2.len());
}
\end{lstlisting}

Spec variables are similar to ghost variables in Dafny or erased values in \fstar.
However, \verus also supports {\it proof variables}, declared as \verb`#[proof]`,
which do not have a correspondence in Dafny or \fstar.
Proof variables sit midway between \verb`exec` variables and \verb`spec` variables:
like \verb`spec` variables, they are ghost and are not compiled to machine code,
but like \verb`exec` variables, they are checked for linearity.
By default, variables in \verb`exec` functions are \verb`exec`,
while variables in \verb`proof` functions are \verb`spec` unless declared \verb`#[proof]`.
Spec functions can only use \verb`spec` variables, not \verb`proof` variables or \verb`exec` variables.

Since \verb`proof` variables are both linear and ghost,
they can represent abstract linear permissions to perform operations,
which can produce and consume the linear permissions.
The next section describes how \verus exploits this feature to verify
safe low-level pointer manipulation that would, in ordinary Rust, require unsafe code.
(In fact, \verus does not support verification of code marked with Rust's \verb`unsafe` keyword;
instead, its goal is to provide safe replacements for unsafe Rust features,
based on linear ghost permissions and SMT-based verification.)




